% !PRINTED 317
\weberpage{317}

In der oben erwähnten Arbeit von \gesperrt{Kiepert} wird die allgemeine Gleichung fünften Grades auf die Form (15) transformiert, wozu nur die Auflösung einer quadratischen Gleichung erforderlich ist. Sieht man darin $a, b, c$ als beliebig gegeben an, so dienen die Gleichungen (16) zur Bestimmung von $\lambda, \mu, \gamma_3$. Dies geschieht ebenfalls mit Hilfe einer quadratischen Gleichung, was allerdings nicht auf den ersten Blick zu ersehen ist. (Vgl. F. \gesperrt{Klein}, Vorlesungen über das Ikosa\"eder, S. 191 f.)

Auf die Resolventenbildung für den siebenten Transformationsgrad werden wir weiter unten zurückkommen, wenn wir über die Hilfsmittel verfügen, die uns die komplexe Multiplikation bietet.

\begin{center}\rule{0.2\linewidth}{0.4pt}\end{center}
